
\subsection{Описание требований к блоку получения исходных данных}

Любая математическая формула и каждая методика расчета, используемая в области естественных наук, независимо от своей сложности и правильности, не может полностью раскрыть весь свой потенциал без стабильного, проверенного и доступного источника точных и полных данных.

В рамках решаемой задачи такими данными будут являться сведения об:
\begin{itemize}
    \item положении береговой линии;
    \item волновых процессах в пирбрежной зоне;
    \item батиметрии;
    % \item рельефе;
    \item климате.
\end{itemize}

Продолжить этот ряд можно геологическими данными, спутниковыми снимками, экономической и топографической информацией...
Продолжать этот ряд можно практически бесконечно.
Но какие из этих данных важнее?
Как их получить и откуда?
Доступны ли они и как долго такими останутся?..
Без ответа на эти вопросы и обеспечения следующего блока расчета необходимыми инструментами получения данных он потеряет весь свой смысл.

Формулируя ранее требования к решению, как универсально применимые для любого участка берега, на источники получения данных накладывается существенное ограничение их \enquote{открытости}.
Таким образом, важность разработки инструментов получения и анализа иходных данных, становится по важности и необходимости наравне с непосредственным блоком расчетов.

С позиции отстраненной от проблем прибрежных территорий, разрабатываемые в работе инструменты сбора данных могут представляться самодостаточно ценными, так как равно эффективно доступны к использованию и для других типов территории вне контекста рассматриваемых задач.

Возвращаясь к форммулированию требований к рассматриваемому блоку методики, естественным и очевидным решением будет наложить на него выполнение задач получения, агрегации, проверки и оценки данных, необходимых для последующего выполнения расчетов.

Многообразие инструментов и разнородность данных требуют для их раскрытия существенного объема описания и сложной рубрикации, в следствие чего частные вопросы по этой теме вынесены в далее следующий раздел~\ref{part:data_intro} на стр.~\pageref{part:data_intro}.

